\subsection{The two proposed quaternion algebras are different}
\paragraph{D23.L7.}
The rational quaternion algebras $(-1,-3)$ and $(2,3)$ are
not isomorphic.  Their ramification sets are, respectively,
$\{3,\infty\}$ and $\{2,3\}$.  Thus the former is definite
and has finite discriminant three; the latter is indefinite
and has discriminant six.  An identification between them cannot
serve as the missing two--three frame comparison.
\begin{proof}
For $(-1,-3)$ the odd-prime Hilbert formula is one away from
three, and at three it is the quadratic residue symbol
$(-1/3)=-1$.  At two the two odd units have residues $7,5$
modulo eight; the exponent in Voight (12.4.14) is
$(({-1}-1)/2)(({-3}-1)/2)=2$, which is even.
The real symbol is negative because both arguments are negative.
This gives its entire ramification set.  The ramification set
of $(2,3)$ was calculated in D23.L3.  A rational algebra
isomorphism would induce an isomorphism over $\mathbf Q_2$,
where the first algebra is split and the second is division,
which is impossible.
\end{proof}
This statement does not dispute a supersingular endomorphism
interpretation of $(-1,-3)$.  It identifies the extra comparison
which would be needed before that interpretation could bear on
the discriminant-six proposal.
